%_____________________________________________________________________________________________
% LATEX Template: Department of Comp/IT BTech Project Reports
% Abstract of Report
% Sun Mar 27 10:34:00 IST 2011
%_____________________________________________________________________________________________
\newpage
\thispagestyle{empty}
\begin{center}
\Large \textbf{Plagiarism Report}
\end{center}
\vspace{1cm}
\noindent [On this page, add the scanned copy of the plagiarism percentage page signed by the Project Guide.]

\newpage
\begin{center}
\Large \textbf{Acknowledgement}
\end{center}
\addcontentsline{toc}{chapter}{Acknowledgement}

We thank our guide, Prof.\ Sunil Bhirud, for his guidance and feedback throughout the project. We thank Dr.\ Pradip Deshmukh, Head of the Department of Computer Science and Engineering, and the faculty and staff of COEP Technological University, Pune, for providing the facilities and support needed to complete this work.

We used open-source tools and public datasets from OpenML, scikit-learn, NumPy, SciPy, pandas, matplotlib and Streamlit, and we thank their maintainers.

\vspace{1cm}
\begin{flushright}
Akshay Naphade\\
Rushikesh Kadepurkar\\
Parth Tatpuje\\
Vishnu Gurav
\end{flushright}

\newpage
\begin{center}
\Large \textbf{Abstract}
\end{center}
\addcontentsline{toc}{chapter}{Abstract}

Machine learning models deployed in production lose accuracy when the data they see drifts away from the data they were trained on. In many applications, such as credit, fraud or churn prediction, the true labels that would reveal this loss arrive days or weeks later. Drift detectors offer earlier, label-free signals, but most input shifts do not harm a given model, and alert thresholds tuned on one dataset often do not hold on another.

This project builds and evaluates an \emph{ML Model Drift Early-Warning System} that tries to warn when the model's accuracy is about to fall, before delayed labels can show it. We developed a benchmark that re-samples real tabular data from OpenML so that inputs drift from a known batch while every row keeps its true label; this lets us measure harm instead of assuming it. The system was developed over ten versions. Versions v1--v5 were exploratory studies on synthetic and natural data streams. Versions v6--v10 were pre-registered: every hypothesis, pass rule and alert threshold was frozen and fingerprinted in code before the confirmation data were downloaded, and each verdict was computed once by a script.

The first pre-registered design (v6), a learned risk calibrator over raw drift scores, failed all five of its hypotheses, because thresholds in dataset-specific units did not transfer: false alerts on unseen datasets ranged from 1.2 to 44.6 per 100 clean batches against a budget of 5. Expressing every signal relative to the stream's own burn-in period fixed this on benchmark streams. The resulting \emph{drift-watch} tier (v7) warned before the accuracy drop on 87\% of harmful gradual shifts on six fresh datasets, a median of six batches ahead, at 4.4 false alerts per 100 clean batches; it failed on natural streams, where drift was already under way when monitoring began. A second \emph{harm} tier (v8) estimates the model's accuracy from an importance-weighted labeled reference set. On five further fresh datasets it cut alarms on harmless shifts from 97\% to 34\% and separated harmful from harmless shifts with an AUROC of 0.78 (95\% CI 0.70--0.86), just below the pre-registered target of 0.80 but clearly above confidence-based baselines (0.56--0.59) on the same data. On a later set of six datasets the same rule exceeded the false-alert budget (12.7 per 100). Two follow-up experiments, nonlinear harm estimators (v9) and stricter alerting rules (v10), did not improve on it and were not adopted.

The project delivers the benchmark and monitoring code, frozen configurations and verdict scripts, a Streamlit dashboard that replays monitored streams batch by batch with delayed labels, a continuous-integration pipeline, and a research paper.

\vspace{0.5cm}
\noindent\textbf{Keywords:} concept drift, dataset shift, model monitoring, delayed labels, performance estimation, importance weighting, pre-registration, MLOps
%_____________________________________________________________________________________________
